%% --- 6.X Native Claim Provenance ---

\subsection{Native Claim Provenance}
\label{sec:exp:native_cpr}

The forensic audit (\S\ref{sec:exp:integrity}) applies uniformly to all systems using only the submitted artifacts.
For \sys{}, which emits structured provenance chains at write-time, we run an additional \emph{native} check that leverages these chains: numerical CPR (Claim Provenance Rate).
This check is unique to \sys{}---no other system in the evaluation produces the required provenance records.

Numerical CPR measures whether quantitative claims in the paper trace to experimental evidence.
During paper generation, the writer annotates each sentence containing a number with a \texttt{\{source: "experimental\_log.md:N"\}} tag linking it to a specific log line.
The claim verifier (\texttt{check\_sources}) extracts the number from the sentence and the referenced log line, then checks whether they match within a 5\% relative tolerance.

Across 15 papers (3 seeds $\times$ 5 tasks), the verifier extracts 639 numerical claims.
Of these, 627 pass (98.1\%).
The 12 failures are predominantly false positives of the extraction heuristic: hardware constants parsed as experimental claims (e.g., ``80GB GPU'' matched against an unrelated log line), LaTeX math subscripts extracted as numbers ($S_{k-1} \to -1.0$), and hyperparameter values described in methodology sections.
Manual inspection finds at most 2--4 genuine mismatches among the 12, yielding a corrected numerical CPR of ${\sim}99\%$.
